% =====================  SECTION II  =====================
\section{Related Work and Research Gap}\label{sec:related}

\subsection{Review Methodology}\label{subsec:review-method}
The literature was surveyed with a focus on works that connect machine learning to
actionable, user-facing financial decisions. Studies were grouped into five
themes---AI-driven wealth management, machine learning for forecasting and distress
prediction, unsupervised risk profiling, explainable AI for financial models, and
decentralized-finance risk together with blockchain auditing---chosen because each maps
onto a stage of the proposed pipeline. Each theme is reviewed for what it contributes
and, more importantly for this work, for the boundary at which it stops. The section
closes with a capability matrix (Table~\ref{tab:capability}) that makes the integration
gap explicit and positions the present contribution against it.

\subsection{AI-Driven Wealth Management and Robo-Advisors}\label{subsec:wealth}
One strand of research examined how artificial intelligence could automate advisory
functions for retail investors, from robo-advisory allocation to predictive analytics
for wealth management \cite{ref1,ref2,ref8}. These systems demonstrated that
personalized, data-driven guidance was feasible at scale and increasingly expected by
users. Multi-agent and collaborative designs pushed this further, coordinating
specialized models for financial research \cite{ref9}. What this line of work
consistently left implicit, however, was the interpretable risk rationale behind an
allocation and any durable record of the advice given; the emphasis was on optimizing
the recommendation, not on making it transparent or verifiable.

\subsection{Machine Learning for Forecasting and Distress Prediction}\label{subsec:forecasting}
Supervised learning was applied to financial forecasting and to predicting financial
distress or bankruptcy, including with multilayer perceptrons \cite{ref3,ref18,ref19};
related work modeled the financial wellbeing of consumers to flag at-risk individuals
\cite{ref7}. These works established that classifiers could identify financial risk
with useful accuracy. At the same time, methodological reviews cautioned that reported
performance in financial machine learning was frequently undermined by leakage and by
causal over-interpretation, and that pitfalls of this kind were common rather than
exceptional \cite{ref4}. This tension---strong headline accuracies alongside a warning
that such accuracies are easily inflated---directly motivates the leakage-aware design
and candid evaluation adopted here.

\subsection{Unsupervised Segmentation and Risk Profiling}\label{subsec:segmentation}
Where labeled risk outcomes were unavailable, unsupervised methods were used to segment
users into behavioral groups; K-means clustering, in particular, was used to profile
investor behavior into interpretable segments \cite{ref20}. This supported the use of
clustering to derive proxy risk bands when ground-truth outcomes were absent. The proxy
nature of such labels is a well-known caveat: a cluster is a descriptive artifact, not a
validated risk outcome, and any downstream classifier trained on it inherits that
limitation---a point this work makes central rather than incidental.

\subsection{Explainable AI for Financial Models}\label{subsec:xai}
As neural models entered financial practice, attribution methods that explained a
prediction in terms of its inputs became a research focus. Benchmarks of
feature-importance methods for neural networks showed that no single method dominated
and that combining methods was often advisable \cite{ref10}, and sensitivity-analysis
techniques were developed specifically for neural networks \cite{ref21}. Explainability
was also applied directly to crypto-asset allocation \cite{ref12}. These methods were
typically studied as standalone techniques; comparatively little work embedded a
per-decision explanation inside a running recommendation pipeline so that every issued
decision carried its own attribution.

\subsection{DeFi Risk, Stablecoins, and Blockchain Audit Trails}\label{subsec:defi}
The remaining literature addressed the decentralized-finance setting and the
verifiability of records. Empirical studies documented liquidation dynamics and
instabilities in DeFi \cite{ref5} and questioned the stability of stablecoins
\cite{ref6}, establishing that crypto allocation carried structural risks beyond price
volatility. Separately, a substantial literature developed blockchain-based,
tamper-evident audit trails and reviewed blockchain applications in financial services
\cite{ref13,ref16,ref22}, including the oracle-trust problem of connecting external data
to on-chain records \cite{ref14,ref15}. Behavioral work added that disclosing a
system's limitations improved user trust and reliance \cite{ref11,ref17}. Each of these
strands secured or motivated one property in isolation; none integrated verifiable
auditing with an explainable, risk-constrained recommendation for an individual user.

\subsection{Comparative Literature Analysis}\label{subsec:comparative}
Table~\ref{tab:capability} summarizes the reviewed themes against the six capabilities
that the proposed system integrates. The pattern is consistent: existing works
implemented one or two capabilities and left the remainder unaddressed. The contribution
of this paper is not the invention of any single column but their coherent unification
in one inspectable system, together with an explicit leakage analysis that few applied
studies report.

\begin{table}[t]
\centering
\caption{Capability comparison of representative prior work and this paper}
\label{tab:capability}
\small
\setlength{\tabcolsep}{4pt}
\begin{tabular}{@{}p{2.6cm} *{6}{>{\centering\arraybackslash}p{1.7cm}}@{}}
\toprule
Representative work & Risk assmt. & XAI & Guardrail & Allocation & Audit log & Leakage anal. \\
\midrule
Robo-advisory / allocation \cite{ref8,ref12} & \partialmark & \partialmark & \nomark & \yes & \nomark & \nomark \\
Distress prediction \cite{ref18,ref19} & \yes & \partialmark & \nomark & \nomark & \nomark & \nomark \\
K-means profiling \cite{ref20} & \partialmark & \nomark & \nomark & \nomark & \nomark & \nomark \\
XAI for finance \cite{ref10,ref21} & \nomark & \yes & \nomark & \nomark & \nomark & \nomark \\
Blockchain audit \cite{ref13,ref22} & \nomark & \nomark & \nomark & \nomark & \yes & \nomark \\
Trust \& disclosure \cite{ref11,ref17} & \nomark & \partialmark & \nomark & \nomark & \nomark & \partialmark \\
This work & \yes & \yes & \yes & \yes & \yes & \yes \\
\bottomrule
\end{tabular}
\par\vspace{3pt}
{\footnotesize\raggedright\textit{Legend:} \yes\ implemented and evaluated;\quad
\partialmark\ discussed or partially addressed;\quad \nomark\ not addressed. The audit
log in this work is an in-memory simulated ledger (Section~\ref{subsec:audit}).\par}
\end{table}

\subsection{Research Gap and Positioning}\label{subsec:positioning}
Read together, the literature established each ingredient of a trustworthy financial
recommender but not their integration. The gap this paper occupies is the assembly of
interpretable risk assessment, per-decision explanation, a conservative safety layer,
risk-aligned allocation, and verifiable logging into one pipeline whose properties are
enforced by construction and whose weaknesses---chiefly the dominance of the rule layer
over the learned model, and the arithmetic leakage between labels and features---are
measured and reported. This positioning deliberately frames the contribution as one of
honest systems integration rather than algorithmic novelty.
